% !TeX encoding = UTF-8
\documentclass[institutional,screen,none]{ua-engineering-v6-beamer}
% !TeX encoding = UTF-8
\UASetUniversity{Universidad Austral}
\UASetFaculty{Facultad de Ingeniería}
\UASetDepartment{Departamento de Computación}
\UASetCampus{Pilar}
\UASetCareer{Ingeniería Informática}
\UASetCommission{Comisión A}
\UASetProfessor{Equipo docente}
\UASetSemester{Segundo cuatrimestre 2026}
\UASetCourse{Sistemas Distribuidos}
\UASetDate{2026}


\UASetClassNumber{02}
\UASetClassTitle{RPC, timeouts, retries e idempotencia}

\title{\UACourse}
\subtitle{Clase 02 --- RPC, timeouts, retries e idempotencia}

\author{\UAProfessor}
\date{\UADate}

\begin{document}

\begin{frame}
\titlepage
\end{frame}

\UASectionSlide{Comunicación en sistemas distribuidos}

\begin{frame}{Objetivos de la clase}
\begin{itemize}
\item Entender una comunicación distribuida como una interacción no inmediata.
\item Razonar qué puede estar ocurriendo mientras el cliente todavía no observa respuesta.
\item Identificar las capas que intervienen en una llamada remota.
\item Analizar cómo cada capa puede detectar, ocultar, transformar o amplificar una falla.
\item Conectar timeout, retry e idempotencia con decisiones de mitigación.
\end{itemize}
\end{frame}

\begin{frame}{Problema real}
En sistemas distribuidos:

\begin{itemize}
\item todo ocurre mediante mensajes
\item la red es incierta
\item no hay garantía de entrega
\end{itemize}

\begin{block}{Pregunta}
¿Qué significa ``la operación ocurrió''?
\end{block}
\end{frame}

\begin{frame}{Comunicar no es instantáneo}
Una operación remota atraviesa varias capas, y cada una puede responder distinto ante una falla.

\begin{center}
\begin{tikzpicture}[
  node distance=0.4cm,
  every node/.style={font=\small},
  box/.style={draw, rounded corners, minimum width=3.0cm, minimum height=0.5cm}
]
\node[box] (app) {Aplicación};
\node[box, below=of app] (rpc) {RPC / cliente HTTP};
\node[box, below=of rpc] (tcp) {Transporte};
\node[box, below=of tcp] (net) {Red};
\node[box, below=of net] (srv) {Servidor};
\draw[->] (app)--(rpc);
\draw[->] (rpc)--(tcp);
\draw[->] (tcp)--(net);
\draw[->] (net)--(srv);
\end{tikzpicture}
\end{center}

\begin{itemize}
\item una capa puede reintentar mientras otra ya venció su timeout
\item una capa puede ocultar pérdida y otra observar demora
\item el servidor puede ejecutar aunque el cliente ya haya abandonado
\end{itemize}
\end{frame}

\begin{frame}{RPC (Remote Procedure Call)}
Abstracción:

\begin{itemize}
\item llamada local: función
\item llamada remota: red
\end{itemize}

\begin{uawarn}
RPC intenta ocultar la existencia de la red. En un sistema robusto se gestiona la red, no se oculta.
\end{uawarn}
\end{frame}

\begin{frame}{Problema clave}
\begin{center}
\Large
¿La operación se ejecutó o no?
\end{center}

\begin{itemize}
\item request perdido
\item response perdido
\item servidor lento
\item servidor caído
\end{itemize}
\end{frame}

\begin{frame}{Timeout}
El cliente decide:

\begin{itemize}
\item esperar
\item reintentar
\item informar error
\end{itemize}

\begin{block}{Clave}
Timeout no implica falla: solo indica que no hubo respuesta dentro de una ventana.
\end{block}
\end{frame}

\begin{frame}{Retry}
Retry introduce:

\begin{itemize}
\item recuperación ante fallas transitorias
\item duplicación potencial
\item presión adicional sobre el sistema
\end{itemize}

Ejemplo:
\begin{itemize}
\item pago duplicado si el primer cobro ocurrió y se perdió la respuesta
\end{itemize}
\end{frame}

\begin{frame}{Idempotencia}
Definición:

\begin{block}{}
Aplicar la operación múltiples veces produce el mismo resultado observable.
\end{block}

Ejemplo:
\begin{itemize}
\item PUT idempotente
\item POST no necesariamente
\item cobro con request id único como idempotencia implementada
\end{itemize}
\end{frame}

\begin{frame}{Semánticas de entrega}
\begin{itemize}
\item at-most-once: evita duplicación, puede perder ejecución
\item at-least-once: evita pérdida, puede duplicar
\item exactly-once: garantía lógica costosa, no magia de la red
\end{itemize}
\end{frame}

\begin{frame}{Trade-off}
\begin{itemize}
\item confiabilidad versus complejidad
\item latencia versus corrección
\item retry versus duplicación
\item idempotencia versus más estado
\end{itemize}
\end{frame}

\begin{frame}{Cierre}
\begin{block}{Idea clave}
Comunicar en sistemas distribuidos es decidir bajo incertidumbre y reconstruir qué pudo haber ocurrido en cada capa.
\end{block}
\end{frame}

\end{document}
